\section{Trabalhos Relacionados}
\label{sec:trabalhos_relacionados}

Os fundamentos do treinamento de redes neurais multicamadas através da descida de gradiente e retropropagação foram consolidados no trabalho seminal de Rumelhart et al. \cite{rumelhart1986learning} e sistematizados de forma abrangente por Goodfellow et al. \cite{goodfellow2016deep}. Desde então, diversos arcabouços computacionais e arquiteturas modernas, incluindo modelos baseados em mecanismos de autoatenção \cite{vaswani2017attention}, foram propostos para escalar o cálculo automático de gradientes e a construção de topologias profundas.

Na literatura contemporânea, o aprendizado de representações e a inferência estruturada demonstram ampla versatilidade em domínios complexos, abrangendo desde a inferência em grafos sobre rastros de rede \cite{nabeel2020following} até a classificação textual com redes neurais e aprendizado de máquina automatizado \cite{muzakir2025text} e a construção de corpora anotados para detecção de padrões em mídias sociais \cite{alvarezojeda2025precom}. Tais avanços reforçam a importância de compreender com rigor analítico e experimental os mecanismos internos de propagação direta e reversa de gradientes investigados neste estudo.
